\documentclass[10pt,a4paper]{amsart}
\usepackage{amsmath,amssymb,amsthm,graphicx}
% Retained preamble for the Stacks Project chapter {CH} (shared project preamble).
% The chapter file pinned-original-source.tex only carries \input{{preamble}} on its first line,
% so no preamble text lives inside the chapter itself.  The Stacks Project uses one single shared
% preamble for every chapter: the snapshot stages it three times (sources/stacks-chow/preamble.tex,
% sources/stacks-etale-cohomology/preamble.tex and focused-seeds/AGP/sources/stacks-snapshot-20261011/
% preamble.tex) and those copies are identical in size.  The block below is copied mechanically,
% line for line, from sources/stacks-chow/preamble.tex lines 14-18 and 154-217: the Xy-pic setup and
% the native theorem declarations together with the mathematical macros.  Nothing was added, removed
% or redefined; the wrapper class and the amsmath/amssymb/amsthm/graphicx base packages are supplied
% by the assembler.
% For commutative diagrams we use Xy-pic
% Chapter: stacks-cohomology.tex -- Cohomology of Sheaves.
\usepackage[all]{xy}

% We use 2cell for 2-commutative diagrams.
\xyoption{2cell}

% Theorem environments.
%
\theoremstyle{plain}
\newtheorem{theorem}[subsection]{Theorem}
\newtheorem{proposition}[subsection]{Proposition}
\newtheorem{lemma}[subsection]{Lemma}

\theoremstyle{definition}
\newtheorem{definition}[subsection]{Definition}
\newtheorem{example}[subsection]{Example}
\newtheorem{exercise}[subsection]{Exercise}
\newtheorem{situation}[subsection]{Situation}

\theoremstyle{remark}
\newtheorem{remark}[subsection]{Remark}
\newtheorem{remarks}[subsection]{Remarks}

\numberwithin{equation}{subsection}

% Macros
%
\def\lim{\mathop{\mathrm{lim}}\nolimits}
\def\colim{\mathop{\mathrm{colim}}\nolimits}
\def\Spec{\mathop{\mathrm{Spec}}}
\def\Hom{\mathop{\mathrm{Hom}}\nolimits}
\def\Ext{\mathop{\mathrm{Ext}}\nolimits}
\def\SheafHom{\mathop{\mathcal{H}\!\mathit{om}}\nolimits}
\def\SheafExt{\mathop{\mathcal{E}\!\mathit{xt}}\nolimits}
\def\Sch{\mathit{Sch}}
\def\Mor{\mathop{\mathrm{Mor}}\nolimits}
\def\Ob{\mathop{\mathrm{Ob}}\nolimits}
\def\Sh{\mathop{\mathit{Sh}}\nolimits}
\def\NL{\mathop{N\!L}\nolimits}
\def\CH{\mathop{\mathrm{CH}}\nolimits}
\def\proetale{{pro\text{-}\acute{e}tale}}
\def\etale{{\acute{e}tale}}
\def\QCoh{\mathit{QCoh}}
\def\Ker{\mathop{\mathrm{Ker}}}
\def\Im{\mathop{\mathrm{Im}}}
\def\Coker{\mathop{\mathrm{Coker}}}
\def\Coim{\mathop{\mathrm{Coim}}}

% Boxtimes
%
\DeclareMathSymbol{\boxtimes}{\mathbin}{AMSa}{"02}

%
% Macros for moduli stacks/spaces
%
\def\QCohstack{\mathcal{QC}\!\mathit{oh}}
\def\Cohstack{\mathcal{C}\!\mathit{oh}}
\def\Spacesstack{\mathcal{S}\!\mathit{paces}}
\def\Quotfunctor{\mathrm{Quot}}
\def\Hilbfunctor{\mathrm{Hilb}}
\def\Curvesstack{\mathcal{C}\!\mathit{urves}}
\def\Polarizedstack{\mathcal{P}\!\mathit{olarized}}
\def\Complexesstack{\mathcal{C}\!\mathit{omplexes}}
% \Pic is the operator that assigns to X its picard group, usage \Pic(X)
% \Picardstack_{X/B} denotes the Picard stack of X over B
% \Picardfunctor_{X/B} denotes the Picard functor of X over B
\def\Pic{\mathop{\mathrm{Pic}}\nolimits}
\def\Picardstack{\mathcal{P}\!\mathit{ic}}
\def\Picardfunctor{\mathrm{Pic}}
\def\Deformationcategory{\mathcal{D}\!\mathit{ef}}
\usepackage{cleveref}
\crefname{originallemma}{original lemma}{original lemmas}
\Crefname{originallemma}{Original lemma}{Original lemmas}
\makeatletter
\newlabel{lemma-cech-map-into}{{01EL}{source}{}{sourcepointer-1}{}}
\newlabel{lemma-cech-map-into@cref}{{[originallemma][0][]01EL}{[1][0][]source}{}{}{}}
\makeatother
\title{The extension-by-zero Cech complex is a resolution of the image presheaf of the augmentation}
\author{The Stacks Project Authors}
\date{stacks-project:cohomology.tex:local-snapshot-20261011; GNU FDL 1.2 or later}
\begin{document}
\maketitle
\noindent\textit{Editorial note.} All delivered mathematical text is the author's own, copied from the pinned Stacks Project chapter file cohomology.tex; no statement or proof step was rewritten, reordered or removed. The author's shorthand (i\_0 = i\_fix) for the insertion homotopy and the eqnarray* computation are reproduced verbatim.\par\medskip
\noindent\textit{Editorial note.} The chapter file pinned-original-source.tex carries only \textbackslash{}input\{preamble\} on its first line, so no preamble text lives inside the chapter. The retained file is a mechanical, line-for-line copy of the project's single shared preamble as staged in the snapshot (the same Stacks-wide preamble of 217 lines that the already published chow and etale-cohomology cards use), namely sources/stacks-chow/preamble.tex lines 14-18 (Xy-pic setup) and 154-217 (theorem declarations and mathematical macros). Nothing was added, removed or redefined, no package was added, and the excerpt serves only as preamble, never as a span coordinate.\par\medskip
\noindent\textit{Editorial note.} The same-chapter pointer is located at its \textbackslash{}label line inside the pinned chapter file.\par\medskip
\noindent\textit{Allowed standard prerequisites (editorial).}\begin{itemize}
\item The extension-by-zero chain complex K(U)\_bullet attached to a covering. Same chapter, Lemma labelled lemma-cech-map-into (source line 1138, Stacks tag 01EL, counted in this batch); named in the counted statement at line 1206. Supplies the complex whose homology presheaves are computed; the computation itself is the counted content.
\item Middle of a finite support argument: sections of the extension-by-zero sheaves. Same chapter, the immediately preceding discussion of the presheaves (j\_\{...\})\_\{p!\}O\_\{...\}, together with the statement of the previous lemma. Supplies the simple description of the presheaves (j\_\{i\_0...i\_p\})\_\{p!\}O\_\{U\_\{i\_0...i\_p\}\} used at proof lines 1243-1244 to rewrite the terms of the extended complex over an open W.
\end{itemize}\medskip
\ifdefined\hypertarget\hypertarget{sourcepointer-1}{}\fi
\noindent\textit{Source reference.} Original lemma 01EL: sources/stacks-cohomology/pinned-original-source.tex, source lines 1138--1138. Same-chapter pointer named in the counted statement at line 1206 for the construction of the complex K(U)\_bullet; the section-wise computation of its homology is retained verbatim in the counted span.\par\medskip
\begin{lemma}
\label{lemma-homology-complex}
Let $X$ be a ringed space.
Let $\mathcal{U} : U = \bigcup_{i \in I} U_i$ be an open covering.
Let $\mathcal{O}_\mathcal{U} \subset \mathcal{O}_X$
be the image presheaf of the map
$\bigoplus j_{p!}\mathcal{O}_{U_i} \to \mathcal{O}_X$.
The chain complex $K(\mathcal{U})_\bullet$ of presheaves
of Lemma \ref{lemma-cech-map-into} above has homology presheaves
$$
H_i(K(\mathcal{U})_\bullet) =
\left\{
\begin{matrix}
0 & \text{if} & i \not = 0 \\
\mathcal{O}_\mathcal{U} & \text{if} & i = 0
\end{matrix}
\right.
$$
\end{lemma}

\begin{proof}
Consider the extended complex $K^{ext}_\bullet$ one gets by putting
$\mathcal{O}_\mathcal{U}$ in degree $-1$ with the obvious map
$K(\mathcal{U})_0 =
\bigoplus_{i_0} (j_{i_0})_{p!}\mathcal{O}_{U_{i_0}} \to
\mathcal{O}_\mathcal{U}$.
It suffices to show that taking sections of this extended complex over
any open $W \subset X$ leads to an acyclic complex.
In fact, we claim that for every $W \subset X$ the complex
$K^{ext}_\bullet(W)$ is homotopy equivalent to the zero complex.
Write $I = I_1 \amalg I_2$ where $W \subset U_i$ if and only
if $i \in I_1$.

\medskip\noindent
If $I_1 = \emptyset$, then the complex $K^{ext}_\bullet(W) = 0$ so there is
nothing to prove.

\medskip\noindent
If $I_1 \not = \emptyset$, then
$\mathcal{O}_\mathcal{U}(W) = \mathcal{O}_X(W)$
and
$$
K^{ext}_p(W) =
\bigoplus\nolimits_{i_0 \ldots i_p \in I_1} \mathcal{O}_X(W).
$$
This is true because of the simple description of the presheaves
$(j_{i_0 \ldots i_p})_{p!}\mathcal{O}_{U_{i_0 \ldots i_p}}$.
Moreover, the differential of the complex $K^{ext}_\bullet(W)$
is given by
$$
d(s)_{i_0 \ldots i_p} =
\sum\nolimits_{j = 0, \ldots, p + 1} \sum\nolimits_{i \in I_1}
(-1)^j s_{i_0 \ldots i_{j - 1} i i_j \ldots i_p}.
$$
The sum is finite as the element $s$ has finite support.
Fix an element $i_{\text{fix}} \in I_1$. Define a map
$$
h : K^{ext}_p(W) \longrightarrow K^{ext}_{p + 1}(W)
$$
by the rule
$$
h(s)_{i_0 \ldots i_{p + 1}} =
\left\{
\begin{matrix}
0 & \text{if} & i_0 \not = i_{\text{fix}} \\
s_{i_1 \ldots i_{p + 1}} & \text{if} & i_0 = i_{\text{fix}}
\end{matrix}
\right.
$$
We will use the shorthand
$h(s)_{i_0 \ldots i_{p + 1}} = (i_0 = i_{\text{fix}}) s_{i_1 \ldots i_p}$
for this. Then we compute
\begin{eqnarray*}
& & (dh + hd)(s)_{i_0 \ldots i_p} \\
& = &
\sum_j \sum_{i \in I_1} (-1)^j h(s)_{i_0 \ldots i_{j - 1} i i_j \ldots i_p}
+
(i = i_0) d(s)_{i_1 \ldots i_p} \\
& = &
s_{i_0 \ldots i_p} +
\sum_{j \geq 1}\sum_{i \in I_1}
(-1)^j (i_0 = i_{\text{fix}}) s_{i_1 \ldots i_{j - 1} i i_j \ldots i_p}
+
(i_0 = i_{\text{fix}}) d(s)_{i_1 \ldots i_p}
\end{eqnarray*}
which is equal to $s_{i_0 \ldots i_p}$ as desired.
\end{proof}

\end{document}
